%----------------------------------------------------------------------------------------------------
% START OF CHAPTER
%----------------------------------------------------------------------------------------------------
\part{Adjectives} \label{M3-Adjectives}

\chapter{Adjectives: Gender and Number Agreement}

Every adjective must agree with the noun in \textbf{gender and number}.

\section{Type 1 Adjectives - ends in -o}
Four forms (dictionary form ends in -o):

\begin{table}[ht]
    \centering
    \begin{tabular}{r c c}
                 & Masculine  & Feminine \\
        \hline
        Singular & \iti{alto} & \iti{alta} \\
        Plural   & \iti{alti} & \iti{alte} \\
    \end{tabular}
\end{table}

\section{Type 2 Adjectives - ends in -e}
Two forms (dictionary form ends in -e):

\begin{table}[ht]
    \centering
    \begin{tabular}{r c c}
        & Masculine  & Feminine \\
        \hline
        Singular & \iti{intelligente} & \iti{intelligente} \\
        Plural   & \iti{intelligenti} & \iti{intelligenti} \\
    \end{tabular}
\end{table}

\section{Adjective Placement}
\begin{itemize}[nosep]
    \item Usually \textbf{after} the noun: \iti{una macchina rossa}
    \item Common adjectives that go \textbf{before}: \iti{bello, brutto, grande, piccolo, giovane, vecchio, buono, cattivo}
\end{itemize}

\section{\itb{Bello} before a noun}
Follows the definite article pattern:
\begin{table}[ht]
    \centering
    \begin{tabular}{r l}
                  & Form \\
        \hline
        \iti{il}  & → \iti{bel} \\
        \iti{lo}  & → \iti{bello} \\
        \iti{la}  & → \iti{bella} \\
        \iti{l'}  & → \iti{bell'} \\
        \iti{i }  & → \iti{bei} \\
        \iti{gli} & → \iti{begli} \\
        \iti{le}  & → \iti{belle} \\
    \end{tabular}
\end{table}

\section{\itb{Buono} before a noun}
Follows the indefinite article pattern:

\iti{buon ragazzo, buono studente, buona ragazza, buon'amica}

\section{Mixed Gender Groups}
When a group contains both genders, adjective defaults to \textbf{masculine plural}:
- \iti{Marco e Sara sono stanchi.}

\chapter{Colours, Nationalities, Personality Traits}

\section{Invariable Colours}
Never change: \iti{rosa, blu, viola, arancione}

\section{Type 1 Colours}
Four forms: \iti{rosso/rossa/rossi/rosse, nero, bianco, giallo, azzurro}

\section{Type 2 Colours}
Two forms: \iti{verde/verdi}

\section{Nationalities}
Not capitalized in Italian.
\begin{itemize}[nosep]
    \item Type 1: \iti{italiano, americano, messicano}
    \item Type 2: \iti{francese, inglese, irlandese, cinese, giapponese}
\end{itemize}

No article is needed after \iti{essere} with nationalities or professions:
 - \iti{Sono italiano. È studentessa.}

\section{Personality Traits}
\begin{itemize}[nosep]
    \item Type 1: \iti{simpatico, antipatico, timido, onesto, pigro, stanco, coraggioso, nervoso}
    \item Type 2: \iti{intelligente, gentile, paziente}
\end{itemize}

\textbf{!!} \iti{Simpatico} = likeable/nice (NOT sympathetic). \iti{Comprensivo} = sympathetic.

\section{-ico Plural Rule}
\begin{itemize}[nosep]
    \item 2-syllable -ico words → add H: \iti{stanco → stanchi, bianco → bianchi}
    \item 3+ syllable -ico words → no H: \iti{simpatico → simpatici, medico → medici}
    \item Feminine plural always adds H before -e: \iti{simpatiche, stanche}
\end{itemize}

\chapter{Possessive Adjectives and Pronouns}

Possessives agree with the \textbf{noun}, not the owner.
\begin{table}[ht]
    \centering
    \begin{tabular}{r c c c c }
            & m. sg. & f. sg. & m. pl. & f. pl. \\
        \hline
        my          & \iti{mio} & \iti{mia} & \iti{miei} & \iti{mie} \\
        your (sg.)  & \iti{tuo} & \iti{tua} & \iti{tuoi} & \iti{tue} \\
        his/her/its & \iti{suo} & \iti{sua} & \iti{suoi} & \iti{sue} \\
        our         & \iti{nostro} & \iti{nostra} & \iti{nostri} & \iti{nostre} \\
        your (pl.)  & \iti{vostro} & \iti{vostra} & \iti{vostri} & \iti{vostre} \\
        their       & \iti{loro} & \iti{loro} & \iti{loro} & \iti{loro} \\
    \end{tabular}
\end{table}

\section{Article Rules}

Always use the definite article before possessives:
\iti{il mio libro, la tua casa, i suoi amici}.

\textbf{!!} Drop the article with singular family members:
\iti{mio padre, mia madre, tuo fratello, sua sorella}

The article \textbf{returns} with:
\begin{itemize}[nosep]
    \item Plural family members: \iti{i miei fratelli}
    \item \iti{Loro: il loro padre}
    \item Adjective present: \iti{il mio vecchio zio}
    \item Diminutives: \iti{il mio fratellino}
\end{itemize}

\section{Suo vs Loro}
\begin{itemize}[nosep]
    \item \iti{suo/sua/suoi/sue} = his/her/its (one owner)
    \item \iti{loro} = their (multiple owners) — \textbf{never changes form}
\end{itemize}

\section{Possessive Pronouns}
Keep the article when the noun is dropped: \iti{Il tuo libro è interessante, ma il mio è noioso.}

%----------------------------------------------------------------------------------------------------
% END OF CHAPTER
%----------------------------------------------------------------------------------------------------